\subsection{vfill and vspace commands}

\begin{frame}{A list without vfill or vspace}

Consider the following list:

\begin{itemize}

\item Item one

\item Item two 

\item Item three 

\item Item four
\end{itemize} 

Without using any commands to space out the items on the list, everything on this slide is cramped together.

\end{frame}

\begin{frame}{A list with vspace}

On this slide, the command vspace is used to manually introduce vertical space between list items, as well as between this top line of text and the list.

\vspace{2em}

\begin{itemize}

\item Item one

\vspace{1em}

\item Item two 

\vspace{1em}

\item Item three 

\vspace{1em}

\item Item four

\vspace{1em}
\end{itemize} 

\end{frame}

\begin{frame}{A list with vfill}

The command vfill provides another option for spacing out text on a slide. Instead of manually controlling the amount of vertical space, this command stretches out items in order to fill the slide.  For example, placing vfill between this top line of text and the list produces the following output:

\vfill
\begin{itemize}

\item Item one

\item Item two 

\item Item three 

\item Item four

\end{itemize} 


\end{frame}




